\chapter{Proposed Methodology}\label{chap:methodology}
This chapter presents the methodology used to construct a simulation-oriented digital twin of a railway network. The 
objective is not to build a complete real-time digital twin, but to define the main steps required to transform 
heterogeneous railway datasets into SUMO-compatible simulation models.

The proposed approach is based on two complementary levels of representation. The macro-level model represents the 
railway network as a station-to-station graph and is mainly used to integrate real operational data into simulation. The 
micro-level model represents the infrastructure in more detail through tracks, switches, platforms, and routing structures. 
This distinction makes it possible to compare a simplified and robust representation with a more detailed but more 
demanding one.

\section{Positioning of This Work}
The state of the art shows that railway digital twins and simulation tools provide useful foundations for analysing 
railway systems. However, constructing an executable simulation from real railway data remains difficult. Infrastructure 
data and operational data are often heterogeneous, incomplete, and represented at different levels of detail. Existing 
simulators also generally assume that the required input files are already available in a consistent and 
simulator-compatible format.

This thesis addresses this gap by focusing on the transformation of railway data into simulation-ready artefacts. The 
contribution is therefore not the development of a new simulator, but the design and evaluation of a data-driven pipeline 
that connects infrastructure reconstruction, operational data integration, route generation, simulation execution, and 
output analysis.

The work is positioned around the comparison of two modelling levels. The macro-level model evaluates how far real 
punctuality data can be integrated into a simplified station-to-station railway simulation. The micro-level model 
investigates whether a more detailed infrastructure representation can support realistic railway routing based on tracks, 
platforms, switches, and a reconstructed dual graph. This comparison is central to the thesis because it highlights the 
trade-off between abstraction, scalability, infrastructure realism, and simulation feasibility.


\section{Framework}
The proposed framework follows a data-driven pipeline that converts raw railway datasets into SUMO-compatible simulation 
scenarios. The pipeline is composed of four main stages: data preprocessing, infrastructure modelling, trip generation, 
and simulation evaluation.

First, the available railway datasets are cleaned and standardized. Infrastructure and operational data must be aligned 
because they do not always share the same identifiers, formats, or level of detail. This step includes filtering the 
selected simulation area, harmonizing station information, processing geographical coordinates, and converting temporal 
values into simulation-compatible timestamps.

Second, the railway infrastructure is reconstructed according to the selected level of abstraction. At the macro level, 
stations are represented as nodes and railway links as edges. This graph is suitable for reconstructing station-to-station 
train movements from punctuality records. At the micro level, the infrastructure is represented through individual track 
segments, detected switches, platform-track associations, and routing structures. This representation is closer to the 
physical infrastructure, but it requires stronger topological consistency.

Third, train movements are generated. In the macro-level model, trips are reconstructed from real punctuality records by 
grouping train passages by train identifier and operating day, then sorting them chronologically. In the micro-level model, 
two strategies are considered: synthetic trip generation and real-data-based trip reconstruction. Synthetic trips are used 
to test network connectivity, while real-data trips aim to reproduce observed train movements on the detailed 
infrastructure.

Finally, the generated networks and routes are executed in SUMO.\@ The simulation outputs are then processed into 
structured datasets that can be analysed and compared with the expected behaviour of the model. This evaluation focuses 
on the quality of the generated infrastructure, the feasibility of train routing, and the ability of each representation 
to support railway simulation.


\section{Why SUMO?}
SUMO was selected because the objective of this thesis requires a flexible, open-source, and scriptable simulation 
platform. Specialized railway simulators such as OpenTrack, RailSys, or OSRD provide strong railway-specific 
functionalities, but they often require detailed and well-structured railway input data. In contrast, this work focuses on 
building a custom pipeline from heterogeneous datasets, which requires full control over network generation, route 
generation, execution, and output processing.

SUMO was selected because it is open-source, scriptable, and compatible with automated data-processing workflows. Its 
XML-based input structure makes it possible to generate networks, routes, stops, and configuration files directly from 
Python scripts. It also provides command-line tools which are useful for converting networks, generating trips, and 
computing executable routes.

SUMO is well suited to this objective because its XML-based structure allows networks, routes, stops, vehicle types, and 
configuration files to be generated automatically from Python scripts. It also provides command-line tools such as 
\texttt{netconvert}, \texttt{randomTrips.py}, and \texttt{duarouter}, which support automated network conversion, 
synthetic trip generation, and route computation.

Another advantage is that SUMO can support both levels of representation considered in this thesis. At the macro level, 
stations can be represented as nodes and station-to-station connections as railway edges. At the micro level, more 
detailed infrastructure can be represented through rail edges, lanes, train stops, and route definitions. This flexibility 
makes SUMO suitable for comparing simplified and detailed railway simulation approaches within the same environment.

However, SUMO is not a railway-only simulator. Detailed signalling, interlocking, dispatching, platform assignment, and 
conflict management are not automatically inferred from raw data. These limitations must be taken into account when 
interpreting the simulations. For this thesis, this trade-off is acceptable because the main objective is to study how 
railway data can be transformed into simulation-ready models, rather than to reproduce all operational rules of a 
professional railway simulator.

\section{Data Available}
The methodology relies on two main categories of data: infrastructure data and operational data. Infrastructure data 
describes the railway network itself, while operational data describes how trains move through this network. Both are 
required to construct a simulation-oriented railway digital twin.

The available data is not directly usable by SUMO.\@ It must first be cleaned, filtered, transformed, and converted into 
simulation files. The type of processing required depends on the level of representation. The macro-level model mainly 
uses stations and station-to-station links, whereas the micro-level model requires track geometry, platforms, switches, 
and routing structures. The data used in this thesis comes from Infrabel Open Data Portal And DTSC who has access to 
private data directly from Infrabel. 

\subsection{Infrastructure Data}
Infrastructure data describes the static railway network. It is used to generate the physical or topological environment 
in which trains circulate. The macro-level and micro-level models use different infrastructure datasets because they 
represent the railway system at different granularities.

At the macro level, the network is represented as a graph of stations and direct connections. This abstraction is 
sufficient for simulating train movements between neighbouring stations, but it does not include platforms, switches, 
signals, or detailed station layouts.

% \subsubsection{Macro-Level Data}

\begin{table}[H]
    \centering
    \begin{tabularx}{\textwidth}{XXXX}
        \toprule
        \textbf{Dataset} & \textbf{Example attributes} & \textbf{Reference value} & \textbf{Role in Methodology} \cr
        \toprule
        Stations & Station identifier, station name, coordinates & 559 stations & Create SUMO nodes \cr
        \midrule
        Stations Connections & Origin, destination, distance, geometry & 1385 tracks & Create SUMO railway edges \cr
        \midrule
        Selected Simulated Area & List of neighbouring stations & User-defined & Restrict the simulation scope \cr
        \bottomrule
    \end{tabularx}
    \caption{Macro-Level Infrastructure Data Overview}\label{tab:macro_level_data}
\end{table}

% \subsubsection{Micro-Level Data}

At the micro level, the infrastructure is represented in more detail. Track segments are used to build the railway graph, 
switches are inferred from the topology, and platforms are associated with nearby track segments. This representation is 
necessary for more realistic routing, but it is more sensitive to missing connections or incorrect topology.

\begin{table}[H]
    \centering
    \begin{tabularx}{\textwidth}{XXXX}
        \toprule
        \textbf{Dataset} & \textbf{Example attributes} & \textbf{Reference value} & \textbf{Role in Methodology} \cr
        \toprule
        Tracks segments & Track identifier, geometry, length & 25068 tracks & Build detailed railway edges \cr
        \midrule
        Stations & Station identifier, station name, coordinates & 554 stations & Support platform detection \cr
        \midrule
        Platforms & Station name, platform number, station's position & 1560 platforms & Create SUMO train stops \cr
        \midrule
        Switches & Node identifier, coordinates, connected tracks & 10379 switches & Support routing and topology 
        analysis \cr
        \bottomrule
    \end{tabularx}
    \caption{Micro-Level Infrastructure Data Overview}\label{tab:micro_level_data}
\end{table}

\subsection{Operational Data}
Operational data describes train movements over the railway network. In this thesis, this information is contained in a 
single punctuality dataset. Each record corresponds to the passage of a train at a station and includes information 
related to train identification, station passage, planned times, observed times, and delays.

This dataset is used to reconstruct ordered train journeys. For each train and operating day, station records are grouped 
and sorted chronologically. At the macro level, these sequences can be mapped onto the station-to-station graph. At the 
micro level, they must be mapped onto a detailed path composed of tracks, platforms, and valid connections, which makes 
the process more complex.

\begin{table}[H]
    \centering
    \begin{tabularx}{\textwidth}{XXXXX}
        \toprule
        \textbf{Dataset} & \textbf{Main attribute groups} & \textbf{Example attributes} & \textbf{Reference value} & 
        \textbf{Role} \cr
        \toprule
        Punctuality dataset & Train identification, station passage, timetable information, observed operation, delays 
        & Train identifier, operating date, station identifier, planned arrival/departure time, observed arrival/departure 
        time, arrival delay, departure delay & Depends on selected time window & Reconstruct real train journeys and 
        generate simulation routes \cr
        \bottomrule
    \end{tabularx}
    \caption{Operational Data Overview}\label{tab:operational_data}
\end{table}

Overall, the available data provides the basis for constructing both macro-level and micro-level railway simulations. The 
infrastructure data defines the network to be simulated, while the operational data defines the train movements to 
reproduce. The central methodological challenge is therefore to align these datasets at the appropriate level of detail and 
convert them into valid SUMO inputs.

% \section{Implementation Plan}